\section{结论与展望}\label{sec:conclusion}

本文围绕机间避碰与全局代数连通保持耦合安全约束下多机器人密集交会的非目标几何死锁、外围静止邻居通道阻挡以及学习残差引起的终点稳态悬停问题，提出了包络等变图策略与邻域近端分配协同控制方法（EGP-NPA）。主要结论如下：
\begin{enumerate}
 \item 通过构造全员可动七机器人对称接触构型的精确有理数二阶 KKT 严格局部极小证书（定理\ref{thm:seven-kkt}），严格证明了构型空间几何与谱安全可达并不蕴含误差单调投影可达，并从代数上导出待互换节点侧向破对称标称驱动与外围静止邻居经由动态配额让渡协同退让的两项充要物理条件。
 \item 提出了满足平面刚体变换群 $\mathrm{SO}(2)\ltimes\mathbb R^2$ 与节点置换群 $S_N$ 等变的一跳消息传递图神经网络残差标称策略（命题\ref{prop:equivariance}），证明了保留 $\mathrm{SO}(2)$ 手性旋转偏置相对于 $\mathrm{O}(2)$ 镜像反射等变在打破对称迎头死锁中的结构必要性（推理期强制 $\mathrm{O}(2)$ 对称化使盲测达标率由 $100\%$ 跌至 $62.5\%$）。
 \item 建立了闭邻域动态分摊护盾下学习探索与跨规模执行的随时硬安全不变性（命题\ref{prop:safety}），导出了接触边界上成对碰撞份额转移驱动静止阻挡邻居同向退让的显式严格正闭式解（命题\ref{prop:unfreeze}），从理论和析因实验两方面证明了静态配额平摊（$B=0$）导致人为冻结、而在损失函数中惩罚法向滤波修正量反而抑制退让驱动力 $\delta_n>0$ 并使达标时间延缓 $+2.00\ \mathrm s$。
 \item 提出了仅依赖一跳目标距离的无状态闭邻域完成包络 $\rho_{i,k}$（命题\ref{prop:envelope}），证明其在保持协同阶段残差幅值的同时于终点邻域将零阶排斥偏置一阶压制为 $\mathcal O(\norm{\be_k}_2)$，消除了密集晶格终点处的极限环悬停；在 $96$ 组训练与跨规模基准及 $48$ 组独立盲测实例上实现了 $100\%$ 达标与零安全违约，并在 \texttt{MuJoCo} 二阶全向底盘闭环中验证了与在线无除法谱估计器的无缝耦合及安全裕度—可通行性权衡定理（定理\ref{thm:passability-tradeoff}）。
 跨规模对比实验表明，所提 EGP-NPA 的安全达标率（$100\%$）显著优于护盾投影一致性安全强化学习（$53.1\%,p=6.1\times 10^{-5}$）、去中心化静态 CBF 护盾图策略（$3.1\%,p=9.3\times 10^{-10}$）、原始—对偶拉格朗日安全强化学习（$0\%$ 无违约达标）、切向环流启发式（$40.6\%,p=3.8\times 10^{-6}$）与短视比例吸引滤波（$15.6\%,p=1.5\times 10^{-8}$）。
\end{enumerate}

在方法边界与未来研究方向上：首先，本文训练分布仅包含 $N\le 10$ 的单层静止邻居包裹晶格，在 $N=25$ 非对称外圈镜像翻转及减半训练预算下仍存在单种子训练方差与手性转向敏感性，未来工作将在节点特征中引入多尺度边界空旷度感知特征并开展大规模多种子课程学习；其次，结合高阶非完整动力学底盘与有向视场角（FOV）感知拓扑的端到端安全协同脱困亦值得深入研究。
